\section*{Question 2}
Two large, finite, parallel conducting plates of area $A$ are placed a distance $d$ apart in vacuum. One plate carries a total charge $+Q$, and the other carries $-Q$. The plates are electrically isolated (i.e., no battery is connected) and no edge effects are considered (i.e., the field is always uniform between the plates). The electric field outside the region between the plates is negligible.

\begin{enumerate}
    \item[(a)] Find the electrostatic energy stored in the system by using the integral 
    \begin{equation*}
        U = \frac{\epsilon_0}{2} \int\limits_{\text{all space}}  E^2 d\tau,
    \end{equation*}
    which is the electrostatic energy stored in an electric field $\vec{E}$, where $\epsilon_0$ is the permittivity of free space and $d\tau$ is a volume element. --- \textit{0.7 pts}
    \item[(b)] Suppose the distance between the plates is increased by a small amount $\Delta x$ with the charges remaining constant. How much work must be done in this process? --- \textit{0.3 pts}
\end{enumerate}

\section*{Solution:}

